\section{Related Work}
\textbf{Tool-Integrated Reasoning and Agentic LLMs.}
Tool use enables LLMs to perform adaptive, multi-step reasoning. Early methods relied on prompt-based~\citep{lu2023chameleon,shen2023hugginggpt} or supervised tool calling~\citep{gou2023tora,qin2023toolllm}, while recent RL-based systems (e.g., RETool~\citep{feng2025retool}, VERL-Tool~\citep{jiang2025verltool}) learn tool-usage policies from environmental feedback, achieving strong performance across QA~\citep{jin2025search}, code-based math reasoning~\citep{xue2025simpletir}, SQL generation~\citep{jiang2025verltool}, and multimodal tasks~\citep{gao2024multi}.
\\
\textbf{Training Collapse in Tool-Integrated GRPO.}
GRPO~\citep{guo2025deepseek} is a value-free, outcome-driven RL method but exhibits severe instability in multi-turn tool-integrated settings, often collapsing when trained from base models with only verifiable rewards~\citep{jin2025search}. This behavior is consistently observed in systems such as Search-R1~\citep{jin2025search}, SimpleTIR~\cite{xue2025simpletir}, and ZeroSearch~\cite{sun2025zerosearch}, whereas PPO remains comparatively stable. Prior explanations attribute the collapse to training–inference mismatch~\citep{liuspeed} or low-likelihood incorrect responses~\citep{xue2025simpletir}, but do not explain PPO’s robustness or the structural origin of the instability. We identify Lazy Likelihood Displacement (LLD) as the underlying mechanism driving collapse in multi-turn TIRL (Section~\ref{sec:d_rela}). Recent works mitigate this via turn-level reward shaping, either relying on ground-truth documents~\cite{zheng2025stepsearch} or an external LLM judge~\citep{zhang2025criticsearch}, both introducing additional supervision or system cost and increased vulnerability to reward hacking~\citep{guo2025deepseek}. Recently, TreeGRPO~\cite{ji2025tree} provides dense reward using tree-based rollout, but increases rollout and computational overhead. In contrast, our approach shows that outcome rewards alone suffice to achieve superior performance once GRPO is stabilized. Detailed related work see~\cref{sec:d_rela}.  
% \xl{It's better to address how the recent methods that address GRPO instability (e.g., SimpleTIR and others) are different from our settings. Otherwise, the reviewer will ask why you did not compare with the related work you mentioned here.}